\paragraph{Statement E4 (Open for the degeneration conclusion; proved
normalization and branch-sign assertions).}\label{res:E4}
\passage{Supplied war-game excerpts of 12 November 2026, entries R4 and H2.
Their references to scans 101--103 have not been independently collated here.}
The normalization assertion H2(a) in the supplied
\texttt{R149\_WARGAME\_EXCERPTS\_20261112.md} is proved in its precise
arithmetic-normalizer formulation: for a hyperbolic curve $C$ over a
number field, a nontrivial closed procyclic subgroup
$I\subset\pi_1(C_{\overline k})$ is essentially cyclotomically
normalized exactly when it is contained in cusp inertia. Thus no such
$I$ exists for a proper hyperbolic curve. Here ``essentially'' means
that a closed subgroup of $N_{\pi_1(C)}(I)$ maps onto an open subgroup
of $G_k$ and acts on $I$ by the cyclotomic character. The same assertion
holds over fields finitely generated over $\mathbf Q$ by AE1.
At a split node with local equation $xy=t$, exchanging the branches
acts on local geometric annulus inertia by inversion. The full
degeneration assertion H2(b), and hence its proposed completion of R4,
is not established in this work and is not used as an input.

\begin{proof}
Stix~\cite[Theorem 20 and \S5.2]{B_Stix2008} proves the stated
equivalence for a characteristic-zero field satisfying
$\operatorname{Hom}_{G_k}(T(S),T(A))=0$ for every torus $S$ and
abelian variety $A$. Part A verifies that hypothesis over number
fields. For the stated finitely generated extension, split $S$ over a
finite extension $k'/k$. A $G_{k'}$-equivariant map from
$T(S)=\widehat{\mathbf Z}(1)^r$ to $T(A)$ has image fixed by
$G_{k'}^\circ$; AE1 makes that image zero. Stix's theorem therefore
applies. In the proper case the set of cusp inertia subgroups is
empty; in the affine case they give exactly the asserted subgroups.
This verifies H2(a) through its actual normalizer map, without choosing
an unproved splitting of that normalizer.

For the branch calculation choose compatible $n$th roots $t_n$ of
$t$ over the geometric generic field, and use the Kummer cover
$z^n=x$. The exchange $x\leftrightarrow y=t/x$ lifts to
$z\mapsto t_n/z$. Conjugating the deck transformation
$z\mapsto\zeta_n z$ by this lift gives
$z\mapsto\zeta_n^{-1}z$. Passing over all $n$ proves the inversion
map on $\widehat{\mathbf Z}(1)$. An arithmetic action additionally
acts on $\zeta_n$ by its cyclotomic character. Thus, when it exchanges
the two branches, its local formula contains that inversion sign as
well. After a finite extension fixing a node and its two branches,
the permutation and branch-exchange issue is removed. This is a
calculation in the local group; it does not assert that the required
global inertia embedding or its arithmetic equivariance has already
been established.
\end{proof}

For a precise record of the remaining source obligation, the opened
Hoshi--Mochizuki manuscript
\cite[Definition 2.10(iv) and Lemma 2.11(ii),(iii)]{B_HoshiMochizuki2011}
gives a local log-group injection whose map on the two node coordinates
is $(a,b)\mapsto a+b$. Write
$I_e^{\log}:=\ker[T^2\xrightarrow{+}T]=\{(a,-a)\}$ for its
geometric node kernel, embedded in the geometric group $\Pi_{G_s}$ of
the associated semigraph. To finish H2(b) in the stated generality one must
supply the log-specialization comparison to the \emph{full} geometric
generic-fibre group and check the displayed maps in the chain
\[
 \widehat{\mathbf Z}(1)
 \xrightarrow{a\mapsto(a,-a)} I_e^{\log}
 \lhook\joinrel\longrightarrow\Pi_{G_s}
 \xrightarrow{\ \mathrm{required\ comparison}\ }
 \pi_1(C_{\overline\eta}),
\]
where the final arrow denotes the required inverse specialization
comparison, not a map constructed in this work. Its obligations include
injectivity, compatibility with the residue-field action,
and the action obtained after \emph{every} splitting of the trait
Galois extension. One must also control the finite permutation of nodes
and branches before asserting individual invariance. The cited local
lemma alone does not provide these entire comparisons. The opened
characteristic-zero Stix theorem also does not by itself prove a
full-group assertion over a finite residue field, where a prime-to-$p$
comparison leaves the $p$-primary part outside its scope.

Accordingly R4(i)'s weight argument and R4(ii)'s proposed passage from
Tate modules to polarized abelian varieties are recorded as unverified
arguments in this deliverable, rather than used as proved inputs. The
smallest retained H2 test is a genus-two degeneration whose special
fibre is two isomorphic CM elliptic curves meeting at one separating
node. Under a surface marking the separating cycle is a commutator
$[a_1,b_1]$, whose image under the actual abelianization map
$\pi\to\pi^{\mathrm{ab}}$ is zero. Thus the abelianized calculation
does not detect that candidate cycle; its global cyclotomic
normalization is exactly the additional assertion requiring the
comparisons above. This status refers to the attached R4 and H2
entries; no independent reading of scans 101--103 is claimed here.
